% =============================================================================
% CONCLUSION
% Best practice: restate purpose, summarize contributions, implications,
% limitations/future work, closing takeaway.
% =============================================================================

\section{Conclusion}
\label{sec:conclusion}

This paper asked why revenue-based proxies mislead research on enterprise digitalization, environmental performance, and productivity, and how widespread that measurement problem is. We did not test whether digitalization improves both TFP and environmental performance (the double dividend). We showed that when researchers use revenue rank for digitalization, ROE for TFP, and revenue growth for environmental performance, regressions capture revenue dynamics rather than the target constructs, and we provided evidence that this pattern is common in the literature.

We documented \emph{prevalence} through a structured audit of 100 firm-level studies: 83\% use at least one revenue-related or accounting-based proxy. We \emph{formalized} same-source proxy bias (Proposition~1), showed that monotone same-source proxies imply a systematic sign and non-vanishing bias (Corollary~1), and that the $t$-statistic can diverge as the sample grows (Theorem~2). We gave a three-step diagnostic and a practical checklist so that applied researchers can assess whether their estimates are likely driven by proxy structure. We \emph{demonstrated} measurement failure with a transparent proxy-based illustration on Chinese A-share firms in heavily polluting industries and with a controlled experiment in which true digitalization is by design uncorrelated with revenue yet the proxy regression remains significant. We also provided a \emph{blueprint} for valid inference: text-based digitalization measures, production-function TFP, emission or GTFP for environmental performance, and strong instruments where needed.

The implications are twofold. For applied researchers, the diagnostic (Section~\ref{sec:formal-diagnostic}) offers a concrete way to check for same-source risk before interpreting proxy regressions as causal. For the literature on digitalization and environment or productivity, our audit suggests that many existing results may conflate revenue or accounting variation with the theoretical constructs; replication and re-estimation with direct measures would clarify the evidence. Limitations are discussed in Section~\ref{sec:discussion} (sample, data merge, audit scope, minimal DGP). Future work could extend the audit to more studies and journal tiers, replicate the diagnostic on published specifications, and test the double dividend with direct measures and quasi-experimental variation.

We do not claim substantive findings on the double dividend. We claim a clear methodological contribution: the measurement problem is common, identification of the target relationship requires direct measures and a design that avoids same-source proxy structure, and the paper supplies a formal characterization, a diagnostic, and a path for rigorous testing in a policy-relevant domain.
